
\documentclass{article}
\usepackage{colortbl}
\usepackage{makecell}
\usepackage{multirow}
\usepackage{supertabular}

\begin{document}

\newcounter{utterance}

\centering \large Interaction Transcript for game `cladder', experiment `full\_v1.5\_default', episode 916 with selfrefine{-}Qwen3.5{-}9B.
\vspace{24pt}

{ \footnotesize  \setcounter{utterance}{1}
\setlength{\tabcolsep}{0pt}
\begin{supertabular}{c@{$\;$}|p{.15\linewidth}@{}p{.15\linewidth}p{.15\linewidth}p{.15\linewidth}p{.15\linewidth}p{.15\linewidth}}
    \# & \multicolumn{2}{c}{Player} && \multicolumn{2}{c}{Game Master} \\
    \hline

    \theutterance \stepcounter{utterance}  
    & & & \multicolumn{4}{p{0.6\linewidth}}{
        \cellcolor[rgb]{0.9,0.9,0.9}{
            \makecell[{{p{\linewidth}}}]{
                \texttt{\tiny{[P1$\langle$GM]}}
                \texttt{You are an expert in causal inference. The following question is not a typical commonsense query, but rather a meticulously designed question created by a professor specializing in causal inference, intended to assess the students' mastery of the course content.} \\
\texttt{Imagine a self{-}contained, hypothetical world with only the following conditions, and without any unmentioned factors or causal relationships: Smoking has a direct effect on effort and college admission. Effort has a direct effect on college admission. For nonsmokers and are lazy, the probability of college admission is 47\%. For nonsmokers and are hard{-}working, the probability of college admission is 88\%. For smokers and are lazy, the probability of college admission is 23\%. For smokers and are hard{-}working, the probability of college admission is 54\%. For nonsmokers, the probability of being hard{-}working is 40\%. For smokers, the probability of being hard{-}working is 61\%. If we disregard the mediation effect through effort, would smoking negatively affect college admission?} \\
\texttt{Start your answer with 'Yes' or 'No', followed by additional reasoning or evidence to support your explanation.} \\
            }
        }
    }
     \\ \\

    \theutterance \stepcounter{utterance}  
    & \multicolumn{4}{p{0.6\linewidth}}{
        \cellcolor[rgb]{0.561,0.737,0.561}{
            \makecell[{{p{\linewidth}}}]{
                \texttt{\tiny{[P1$\rangle$GM]}}
                \texttt{No} \\
\\ 
\texttt{**Reasoning:**} \\
\\ 
\texttt{To determine if smoking negatively affects college admission when disregarding the mediation effect through effort, we must calculate the **Direct Effect** of smoking on admission. The direct effect isolates the impact of smoking by holding the mediator (effort) constant, effectively comparing the outcome of smoking versus not smoking for individuals with the same level of effort.} \\
\\ 
\texttt{Mathematically, the direct effect \$DE\$ is defined as:} \\
\texttt{\$\$DE = P(Y\_\{smoker\} = 1 | E=1) {-} P(Y\_\{nonsmoker\} = 1 | E=1) + P(Y\_\{smoker\} = 1 | E=0) {-} P(Y\_\{nonsmoker\} = 1 | E=0)\$\$} \\
\texttt{Alternatively, this can be viewed as the difference in admission probabilities for hard{-}working individuals plus the difference for lazy individuals.} \\
\\ 
\texttt{**Step 1: Calculate the direct effect for Hard{-}Working individuals (\$E=1\$)**} \\
\texttt{*   Admission probability given Smoker \& Hard{-}working: \$P(Y=1|S=1, E=1) = 0.54\$} \\
\texttt{*   Admission probability given Nonsmoker \& Hard{-}working: \$P(Y=1|S=0, E=1) = 0.88\$} \\
\texttt{*   Difference: \$0.54 {-} 0.88 = {-}0.34\$} \\
\\ 
\texttt{**Step 2: Calculate the direct effect for Lazy individuals (\$E=0\$)**} \\
\texttt{*   Admission probability given Smoker \& Lazy: \$P(Y=1|S=1, E=0) = 0.23\$} \\
\texttt{*   Admission probability given Nonsmoker \& Lazy: \$P(Y=1|S=0, E=0) = 0.47\$} \\
\texttt{*   Difference: \$0.23 {-} 0.47 = {-}0.24\$} \\
\\ 
\texttt{**Step 3: Calculate the Average Direct Effect**} \\
\texttt{Since the question asks if smoking negatively affects admission in general (averaging over the population's effort distribution), we must weight these differences by the probability of being hard{-}working or lazy in the source population (or simply observe the sign of the differences if the direction is consistent).} \\
\texttt{*   For hard{-}workers, smoking reduces admission probability by 34 percentage points.} \\
\texttt{*   For lazy individuals, smoking reduces admission probability by 24 percentage points.} \\
\\ 
\texttt{In both subgroups defined by effort, the probability of admission is strictly lower for smokers than for nonsmokers. Therefore, disregarding the mediation path (the change in effort caused by smoking), the direct causal effect of smoking on college admission is negative.} \\
\\ 
\texttt{**Conclusion:** Yes, smoking negatively affects college admission when disregarding the mediation effect through effort.} \\
            }
        }
    }
    & & \\ \\

    \theutterance \stepcounter{utterance}  
    & & & \multicolumn{2}{p{0.3\linewidth}}{
        \cellcolor[rgb]{0.9,0.9,0.9}{
            \makecell[{{p{\linewidth}}}]{
                \texttt{\tiny{[GM$|$GM]}}
                \texttt{No} \\
            }
        }
    }
    & & \\ \\

    \theutterance \stepcounter{utterance}  
    & & & \multicolumn{2}{p{0.3\linewidth}}{
        \cellcolor[rgb]{0.9,0.9,0.9}{
            \makecell[{{p{\linewidth}}}]{
                \texttt{\tiny{[GM$|$GM]}}
                \texttt{yes} \\
            }
        }
    }
    & & \\ \\

    \theutterance \stepcounter{utterance}  
    & & & \multicolumn{2}{p{0.3\linewidth}}{
        \cellcolor[rgb]{0.9,0.9,0.9}{
            \makecell[{{p{\linewidth}}}]{
                \texttt{\tiny{[GM$|$GM]}}
                \texttt{game\_result = LOSE} \\
            }
        }
    }
    & & \\ \\

\end{supertabular}
}

\end{document}
